\documentclass[11pt, letterpaper]{article}
\usepackage[utf8]{inputenc}
\usepackage{geometry}
\geometry{margin=1in}
\usepackage{booktabs}
\usepackage{amsmath}
\usepackage{amssymb}
\usepackage{hyperref}
\usepackage{longtable}
\usepackage{xcolor}
\usepackage{array}
\usepackage{float}
\usepackage{colortbl}

% Define custom colors for clarity
\definecolor{headerblue}{RGB}{0, 51, 102}
\definecolor{rowgray}{RGB}{240, 240, 240}
\definecolor{highlight}{RGB}{255, 255, 204}

\title{\textbf{Comprehensive Guide to 3-Statement Financial Modeling}\\ \large Wall Street Knights Training - Week 2}
\author{Financial Modeling Reference Guide}
\date{\today}

\begin{document}

\maketitle
\tableofcontents
\newpage

\section{Introduction to the 3-Statement Model}

\subsection{What is this model?}
You have been assigned to build a \textbf{3-Statement Model}. This is the fundamental tool used in investment banking, private equity, and corporate finance. It links the three core financial reports into one dynamic, interconnected system:

\begin{enumerate}
    \item \textbf{Income Statement (IS):} Shows profitability over a period of time (Revenue $-$ Expenses $=$ Net Income).
    \item \textbf{Balance Sheet (BS):} Shows the company's financial position at a specific point in time (Assets $=$ Liabilities $+$ Equity).
    \item \textbf{Cash Flow Statement (CFS):} Shows the actual movement of cash in and out of the business, reconciling Net Income to the actual change in the bank account.
\end{enumerate}

\subsection{Why do we build this?}
The primary purpose is to determine if a company is \textbf{viable}.
\begin{itemize}
    \item \textbf{Profit $\neq$ Cash:} A company can show a profit on the Income Statement (Net Income) but go bankrupt because it ran out of cash to pay bills (Liquidity Crisis).
    \item \textbf{Forecasting:} By projecting these statements forward (2021E--2023E), we can predict if the company will need to raise more debt, if it can afford to pay dividends, or if it has enough cash to buy new equipment.
\end{itemize}

\section{Glossary of Abbreviations}
Below are the abbreviations used in your assignment and common in the industry.

\begin{longtable}{@{}p{0.15\textwidth} p{0.35\textwidth} p{0.45\textwidth}@{}}
\toprule
\textbf{Abbr.} & \textbf{Full Term} & \textbf{Definition} \\
\midrule
\endhead

\textbf{IS} & Income Statement & Financial statement showing revenue and expenses. \\
\textbf{BS} & Balance Sheet & Financial statement showing assets, liabilities, and equity. \\
\textbf{CFS} & Cash Flow Statement & Financial statement showing cash inflows and outflows. \\
\textbf{FY / A / E} & Fiscal Year / Actual / Estimated & 2020A means "Actual 2020 results"; 2021E means "Estimated 2021 results." \\
\textbf{COGS} & Cost of Goods Sold & Direct costs attributable to the production of goods sold. \\
\textbf{GP} & Gross Profit & Revenue minus COGS. \\
\textbf{EBITDA} & Earnings Before Interest, Taxes, Depreciation, \& Amortization & A proxy for operating cash flow and core profitability. \\
\textbf{D\&A} & Depreciation \& Amortization & Allocation of the cost of an asset over its useful life (Non-cash expense). \\
\textbf{EBIT} & Earnings Before Interest \& Taxes & Also known as Operating Profit/Income. \\
\textbf{NI} & Net Income & The "Bottom Line" profit after all expenses, interest, and taxes. \\
\textbf{EPS} & Earnings Per Share & Net Income divided by shares outstanding. \\
\textbf{PP\&E} & Property, Plant, \& Equipment & Long-term tangible assets (machinery, buildings). \\
\textbf{CapEx} & Capital Expenditures & Funds used to acquire, upgrade, or maintain physical assets (PP\&E). \\
\textbf{NWC} & Net Working Capital & Current Assets minus Current Liabilities. Measures short-term liquidity. \\
\textbf{AR} & Accounts Receivable & Money owed to the company by customers (Sales made on credit). \\
\textbf{AP} & Accounts Payable & Money the company owes to suppliers (Bills not yet paid). \\
\textbf{Inv} & Inventory & Goods available for sale or raw materials. \\
\textbf{DSO} & Days Sales Outstanding & How many days it takes to collect cash from customers. \\
\textbf{DIO} & Days Inventory Outstanding & How many days inventory sits before being sold. \\
\textbf{DPO} & Days Payable Outstanding & How many days the company takes to pay suppliers. \\
\textbf{Ret. Earnings} & Retained Earnings & Cumulative Net Income minus Dividends paid out. \\
\bottomrule
\end{longtable}

\newpage

\section{The Core Mechanics: How Data Moves}

The model works like a circulatory system. If one "vein" is cut, the Balance Sheet will not balance.

\subsection{The Circular Flow}
\begin{enumerate}
    \item \textbf{Start at Income Statement:} Calculate Revenue down to \textbf{Net Income}.
    \item \textbf{Flow to Cash Flow Statement:} \textbf{Net Income} is the starting point (top line) of the CFS.
    \item \textbf{Flow to Balance Sheet (Equity):} \textbf{Net Income} also increases \textbf{Retained Earnings} in the Equity section of the BS.
    \item \textbf{Adjust for Non-Cash Items (CFS):} Add back Depreciation (from IS/PP\&E Schedule) to the CFS because it lowered Net Income but didn't cost cash.
    \item \textbf{Working Capital (BS $\to$ CFS):} Calculate changes in AR, Inventory, and AP.
    \begin{itemize}
        \item Asset increases (AR, Inv) \textbf{consumes} cash (negative on CFS).
        \item Liability increases (AP) \textbf{provides} cash (positive on CFS).
    \end{itemize}
    \item \textbf{Investing \& Financing (Schedules $\to$ CFS):} 
    \begin{itemize}
        \item CapEx flows from the PP\&E Schedule to Investing Activities.
        \item Debt Repayment flows from the Debt Schedule to Financing Activities.
    \end{itemize}
    \item \textbf{The Anchor (CFS $\to$ BS):} The bottom of the CFS calculates the \textbf{Net Change in Cash}.
    \[ \text{Ending Cash} = \text{Beginning Cash} + \text{Net Change in Cash} \]
    This \textbf{Ending Cash} links to the top line of the Balance Sheet.
\end{enumerate}

\subsection{The Balancing Equation}
If linked correctly:
\[ \text{Total Assets} = \text{Total Liabilities} + \text{Shareholders' Equity} \]
Your model includes a "Check" line. This is calculated as:
\[ \text{Check} = \text{Assets} - (\text{Liabilities} + \text{Equity}) \]
This must always be \textbf{ZERO}.

\newpage

\section{Step-by-Step Assignment Walkthrough}

\subsection{Step 1: Assumptions \& Income Statement (PDF Page 2)}
You must forecast 2021E, 2022E, and 2023E.

\textbf{1. Revenue Growth:}
The assumption table provides growth rates (35.0\%, 25.0\%, 18.0\%).
\begin{equation}
    \text{Revenue}_{2021} = \text{Revenue}_{2020} \times (1 + 35\%)
\end{equation}

\textbf{2. Cost of Goods Sold (COGS):}
The assumption table provides Gross Margin \% (41\%, 42\%, 43\%). You must back-solve for COGS.
\begin{equation}
    \text{COGS} = \text{Revenue} - (\text{Revenue} \times \text{Gross Margin \%})
\end{equation}
\textit{Alternatively: COGS = Revenue * (1 - Gross Margin \%)}

\textbf{3. Operating Expenses (SG\&A / Marketing):}
Usually forecasted as a \% of Revenue. Based on historicals (2018-2020), calculate the \% and apply it to projected revenue.

\textbf{4. Depreciation:}
\textbf{DO NOT} forecast this as a simple \% of sales on the Income Statement. This must link from the \textbf{PP\&E Schedule}.

\textbf{5. Interest Expense:}
This links from the \textbf{Debt Schedule}.
\begin{equation}
    \text{Interest Exp} = \text{Average Debt Balance} \times \text{Interest Rate (6\%)}
\end{equation}

\textbf{6. Taxes:}
\begin{equation}
    \text{Taxes} = (\text{Operating Profit} - \text{Interest Expense}) \times \text{Tax Rate}
\end{equation}

\subsection{Step 2: Balance Sheet Working Capital (PDF Page 3)}
You are given "Days" assumptions: DSO (48), DIO (22), DPO (17).

\textbf{1. Accounts Receivable (AR):}
\begin{equation}
    \text{AR} = \frac{\text{Revenue}}{365} \times \text{DSO (48)}
\end{equation}

\textbf{2. Inventory:}
Inventory relates to COGS, not Revenue.
\begin{equation}
    \text{Inventory} = \frac{\text{COGS}}{365} \times \text{DIO (22)}
\end{equation}

\textbf{3. Accounts Payable (AP):}
Payables relate to COGS (suppliers).
\begin{equation}
    \text{AP} = \frac{\text{COGS}}{365} \times \text{DPO (17)}
\end{equation}

\subsection{Step 3: The Supporting Schedules (PDF Pages 4-5)}

\textbf{A. PP\&E Schedule:}
\begin{itemize}
    \item \textbf{Beg Balance:} Link from Prior Year Ending Balance.
    \item \textbf{Plus CapEx:} Calculated as \% of Revenue (0.9\% provided).
    \item \textbf{Less Depreciation:} Calculated as \% of Revenue (0.4\% provided).
    \item \textbf{End Balance:} Beg + CapEx - Depreciation. (Links to Balance Sheet).
\end{itemize}

\textbf{B. Debt Schedule:}
\begin{itemize}
    \item \textbf{Beg Balance:} Link from Prior Year.
    \item \textbf{Less Repayment:} Hardcoded in PDF (\$1,383 for 2021, 2022, 2023).
    \item \textbf{End Balance:} Beg - Repayment. (Links to Balance Sheet).
\end{itemize}

\textbf{C. Equity Schedule:}
\begin{itemize}
    \item \textbf{Plus Net Income:} Links from Income Statement.
    \item \textbf{Less Dividends:} Hardcoded in PDF (\$8,000).
    \item \textbf{End Balance:} Beg + Net Income - Dividends. (Links to Balance Sheet).
\end{itemize}

\newpage

\section{Visualizing the Flow: The Financial Statements}
Below is a reconstruction of the model structure based on your homework. Arrows ($\rightarrow$) indicate where data flows.

\subsection{Income Statement (Simplified)}
\begin{table}[H]
\centering
\begin{tabular}{lrrrrr}
\toprule
\textbf{($\$$ in Thousands)} & \textbf{2019A} & \textbf{2020A} & \textbf{2021E} & \textbf{2022E} & \textbf{2023E} \\
\midrule
Total Revenue & 33,312 & 52,072 & \textit{Calc} & \textit{Calc} & \textit{Calc} \\
Cost of Goods Sold & (20,652) & (31,309) & \textit{Calc} & \textit{Calc} & \textit{Calc} \\
\textbf{Gross Profit} & \textbf{12,660} & \textbf{20,762} & \textbf{Calc} & \textbf{Calc} & \textbf{Calc} \\
\dots & \dots & \dots & \dots & \dots & \dots \\
Depreciation & (302) & (343) & $\leftarrow$ \textit{Link PP\&E} & \dots & \dots \\
\textbf{Operating Profit} & \textbf{9,055} & \textbf{17,203} & \textbf{Calc} & \dots & \dots \\
Interest Expense & (380) & (249) & $\leftarrow$ \textit{Link Debt} & \dots & \dots \\
Taxes & (2,228) & (4,352) & \textit{Calc} & \dots & \dots \\
\midrule
\textbf{Net Income} & \textbf{6,448} & \textbf{12,602} & \textbf{Calc} & \textbf{Calc} & \textbf{Calc} \\
\bottomrule
\end{tabular}
\end{table}

\subsection{Cash Flow Statement (The Bridge)}
\begin{table}[H]
\centering
\begin{tabular}{lrrr}
\toprule
\textbf{($\$$ in Thousands)} & \textbf{2021E} & \textbf{Link Source} \\
\midrule
\textbf{Cash Flow from Operations (CFO)} & & \\
Net Income & \textit{Value} & $\leftarrow$ \textit{From Income Statement} \\
(+) Depreciation & \textit{Value} & $\leftarrow$ \textit{From PP\&E Sched OR IS} \\
\textit{Changes in Working Capital:} & & \\
(Increase)/Decrease in AR & \textit{(Value)} & $\leftarrow$ \textit{Prior Year AR - Current AR} \\
(Increase)/Decrease in Inv & \textit{(Value)} & $\leftarrow$ \textit{Prior Year Inv - Current Inv} \\
Increase/(Decrease) in AP & \textit{Value} & $\leftarrow$ \textit{Current AP - Prior Year AP} \\
\textbf{Total CFO} & \textbf{Sum} & \\
\midrule
\textbf{Cash Flow from Investing (CFI)} & & \\
Capital Expenditures (CapEx) & \textit{(Value)} & $\leftarrow$ \textit{From PP\&E Sched} \\
\textbf{Total CFI} & \textbf{Sum} & \\
\midrule
\textbf{Cash Flow from Financing (CFF)} & & \\
Debt Repayment & \textit{(1,383)} & $\leftarrow$ \textit{From Debt Sched} \\
Dividends Paid & \textit{(8,000)} & $\leftarrow$ \textit{From Equity Sched} \\
\textbf{Total CFF} & \textbf{Sum} & \\
\midrule
\textbf{Net Change in Cash} & \textbf{Sum} & (CFO + CFI + CFF) \\
Beginning Cash Balance & 6,372 & $\leftarrow$ \textit{2020 Ending Cash} \\
\textbf{Ending Cash Balance} & \textbf{Calc} & $\rightarrow$ \textit{Links to BS Cash} \\
\bottomrule
\end{tabular}
\end{table}

\subsection{Balance Sheet (The Destination)}
\begin{table}[H]
\centering
\begin{tabular}{lrrrr}
\toprule
\textbf{($\$$ in Thousands)} & \textbf{2020A} & \textbf{2021E} & \textbf{Link Source} \\
\midrule
\textbf{Assets} & & & \\
Cash & 6,372 & \textit{Calc} & $\leftarrow$ \textit{From CFS Bottom} \\
Accounts Receivable & 6,894 & \textit{Calc} & $\leftarrow$ \textit{Calc using Days} \\
Inventory & 1,875 & \textit{Calc} & $\leftarrow$ \textit{Calc using Days} \\
PP\&E (Net) & 1,974 & \textit{Calc} & $\leftarrow$ \textit{From PP\&E Sched} \\
\textbf{Total Assets} & \textbf{17,640} & \textbf{Sum} & \\
\midrule
\textbf{Liabilities \& Equity} & & & \\
Accounts Payable & 1,462 & \textit{Calc} & $\leftarrow$ \textit{Calc using Days} \\
Long Term Debt & 4,150 & \textit{Calc} & $\leftarrow$ \textit{From Debt Sched} \\
Shareholders' Equity & 12,028 & \textit{Calc} & $\leftarrow$ \textit{From Equity Sched} \\
\textbf{Total L\&E} & \textbf{17,640} & \textbf{Sum} & \\
\midrule
\textbf{Check} & \textbf{0} & \textbf{0} & (Assets - L\&E) \\
\bottomrule
\end{tabular}
\end{table}

\newpage

\section{Common Questions \& Troubleshooting}

\subsection{Q: Why is my Balance Sheet check not zero?}
This is the most common error. Check the following:
\begin{enumerate}
    \item \textbf{Cash Link:} Did you link the calculated "Ending Cash" from the CFS to the "Cash" line on the Balance Sheet? Do not hardcode it or plug it.
    \item \textbf{Equity Link:} Did \textbf{Net Income} flow into the Equity Schedule (Retained Earnings)?
    \item \textbf{Sign Errors (Working Capital):} On the CFS:
    \begin{itemize}
        \item If AR goes UP from 2020 to 2021, cash flow is \textbf{NEGATIVE}.
        \item If AP goes UP from 2020 to 2021, cash flow is \textbf{POSITIVE}.
    \end{itemize}
    \item \textbf{Schedule Math:} Ensure your Debt Schedule calculates: $Beg - Repayment = End$. Ensure PP\&E calculates: $Beg + CapEx - Depr = End$.
\end{enumerate}

\subsection{Q: What is the "Circ Toggle"?}
"Circ" stands for \textbf{Circularity}.
\begin{itemize}
    \item \textbf{The Loop:} Interest Expense reduces Net Income $\rightarrow$ Reduced Net Income reduces Cash $\rightarrow$ Less Cash means more Debt (in complex models with a Revolver) $\rightarrow$ More Debt means higher Interest Expense.
    \item \textbf{The Toggle:} A switch (usually an IF statement) that breaks this loop calculation to stop Excel from crashing (showing \#REF!).
    \item \textbf{In this Homework:} Since Debt Repayment is hardcoded (\$1,383) and not based on available cash, you likely do not have active circularity, so the toggle is just for show or future complexity.
\end{itemize}

\subsection{Q: Why do we add back Depreciation on the CFS?}
Depreciation is an expense on the Income Statement that reduces Net Income to save on taxes. However, you did not actually write a check for "Depreciation" this year. The cash left the building years ago when you bought the machine (CapEx). Therefore, to get back to "Actual Cash," you must add that paper expense back.

\subsection{Q: How do I handle the "Check" line?}
The check line exists to verify the Accounting Equation ($A = L + E$).
\begin{itemize}
    \item \textbf{Balance Sheet Check:} $Total Assets - (Total Liabilities + Total Equity)$. Result should be 0.
    \item \textbf{Cash Flow Check:} $(Ending Cash Calculated on CFS) - (Cash Balance on BS)$. Result should be 0.
\end{itemize}

\end{document}